% migrate-to: execution-market-making, market-making-lob
\subsection{Queue Dynamics and Latency}

\begin{principle} \hlt{Condition on the Queue}\\
Queue dynamics are less random once the state of the book is specified. A cancellation at the best bid is not comparable when that queue has one lot, ten lots, or a much larger opposite queue.
\end{principle}

\begin{definition} \hlt{Queue Imbalance}\\
Queue imbalance compares available quantity on the best bid and ask. Predictive because a small best queue is easier to deplete, while a large opposite queue can make the next price move less likely to reverse immediately.
\end{definition}

\begin{mechanism} \hlt{Arrival-Departure Ratio}\\
For a queue, the arrival-departure ratio compares insertion intensity with cancellation plus transaction intensity. Ratios above one imply queue growth; ratios below one imply depletion.
\end{mechanism}

\begin{caveat} \hlt{Adverse Queue Selection}\\
A limit order is most likely to fill when marketable flow reaches it. That event is not neutral: fills are often more valuable to the taker precisely when the resting quote is stale, underpriced, or exposed to informed flow.
\end{caveat}

\begin{process} \hlt{Queue-Reactive Estimation}
\begin{enumerate}[label=\roman*.]
\setlength{\itemsep}{0pt}
\item choose a starting time and record the queue state;
\item wait for the next insertion, cancellation, or transaction;
\item store the waiting time by event type and queue-size bucket;
\item estimate state-dependent intensities from the average waiting times.
\end{enumerate}
\end{process}

\begin{caveat} \hlt{Mechanical Mean Reversion}\\
Queue-reactive simulations can generate too much price mean reversion: after a best queue disappears, the newly promoted queue may be large enough to pull prices back. Models often need memory, Hawkes-style excitation, or state resets after price moves.
\end{caveat}

\begin{principle} \hlt{Latency Race Boundary}\\
Latency matters most when a trader can change queue position before others react. It matters less when the tick is wide, the queue is deep, or the order is deliberately not competing for first priority.
\end{principle}
